\section{Pure state quantum information}

Finite-dimensional quantum information theory takes place inside a complex vector space $\mathbb{C}^{d^n}$
and $n$-qudit states are equivalence classes of vectors in this space.
In the continuous-variable setting, just to get as far as saying what a state is takes us on a bit of a journey.
We'll give a whistle-stop tour here, but since the technical details won't be relevant for this thesis, we'll skip all the scary bits.

\subsection{The Hilbert space $\mathcal{L}^2(\mathbb{R}^n)$}

The continuous-variable analogue of a qudit is called a \textit{qumode}.
Throughout, we reserve the right to get sloppy and just call these \textit{modes}.
To define $n$-qumode states, we start with a particular infinite-dimensional complex vector space denoted $\mathcal{L}^2(\mathbb{R}^n)$.
Elements of this space are equivalence classes of certain functions $\mathbb{R}^n \to \mathbb{C}$.
These functions $f$ must be absolutely square-integrable, i.e.\ $\int_{\mathbb{R}^n} dx |f(x)|^2$ must converge.
Strictly speaking, the integral here is the Lebesgue integral rather than the Riemann integral, but this is the first of many technicalities we'll sweep under the carpet.
The equivalence relation is that $f \sim g \iff f - g$ is a null function, where again we won't define null functions -- for us it'll be fine to forget elements are equivalence classes, and just view them as functions.

This space comes with an inner product $\langle f, g \rangle = \int_{\mathbb{R}^n} d\boldsymbol{x} \overline{f(\boldsymbol{x})} g(\boldsymbol{x})$, where the overline denotes complex conjugation.
A \textit{Hilbert space} is defined to be a \textit{complete} inner product space.
It won't be relevant for us to say what this means; all we'll note is that $\mathcal{L}^2(\mathbb{R}^n)$ with the inner product above is indeed complete, so is a Hilbert space.
Furthermore, one can define a \textit{topology} on this space; this is yet another detail we won't go into, but it makes it a \textit{topological vector space}, which lets us talk about notions like \textit{continuity} in the coming subsections.

\teague{Something about this space being decomposable into tensor products of $\LLL^2(\RR)$}

\subsection{Schwartz space and rigged Hilbert spaces}

Within our Hilbert space $\mathcal{H} = \mathcal{L}^2(\mathbb{R}^n)$, there lives a subspace $\Phi \leq \mathcal{H}$ called \textit{Schwartz space}.
This is the space of all functions that are \textit{infinitely differentiable} and \textit{rapidly decaying}.
Intuitively, this means such functions $\psi: \mathbb{R}^n \to \mathbb{C}$ are smooth, and as $|| \boldsymbol{x} || \to \infty$ they quickly become flat and tend to $0$.
This subspace inherits many of the nice properties of $\mathcal{L}^2(\mathbb{R}^n)$; in particular, it too is a topological Hilbert space.

Now, for any topological complex vector space $V$, we can define the continuous dual space $V^*$ and continuous conjugate-dual space $V^\times$.
The continuous dual space is the space of all continuous linear maps from $V$ to the underlying field $\mathbb{C}$.
That is, for $F \in V^*$, the linear part means we have $F(\alpha u + \beta v) = \alpha F(u) + \beta F(v)$.
We won't define what it means for this function to be continuous.
Analogously, for $G \in V^\times$, we have $G(\alpha u + \beta v) = \overline{\alpha} G(u) + \overline{\beta} G(v)$.
We will be interested in the continuous dual and conjugate-dual spaces of Schwartz space $\Phi$ and the Hilbert space $\mathcal{H} = \mathcal{L}^2(\mathbb{R}^n)$ it sits in.
These form two chains of relations:
\begin{equation}\label{eq:rigged_spaces_separate}
\Phi \leq \HHH \cong \HHH^* \leq \Phi^* \quad \text{ and } \quad \Phi \leq \HHH \cong \HHH^\times \leq \Phi^{\times}
\end{equation}
We can squish all this into one diagram, and combine this with the fact that the continuous dual and conjugate-dual of Schwartz space are isomorphic:
\begin{equation}\label{eq:rigged_spaces_combined}
TODO
\end{equation}
Specifically, the isomorphism $\Phi^* \cong \Phi^\times$ can be written $F \mapsto \overline{F}$, where $\overline{F}(\psi) \coloneqq \overline{F(\psi)}$.

\subsection{Pure states}

With that, we're done; we can now define pure qumode states.
We'll divide these into two classes; \textit{physical} and \textit{generalised}.
Physical pure $n$-qumode states are equivalence classes of vectors in the Hilbert space $\mathcal{H} = \mathcal{L}^2(\mathbb{R}^n)$, where the equivalence relation is $f \sim g \iff f = \lambda g$ for some non-zero $\lambda \in \mathbb{C}$.
Generalised states are elements of $\Phi^\times$, i.e.\ continuous conjugate-linear functionals on Schwartz space.
So by \Cref{eq:rigged_spaces_separate}, physical states are a subspace of the generalised states, up to isomorphism.
Indeed, given physical state $f \in \mathcal{H}$, we can identify it with the functional:
\begin{equation}
    \psi \mapsto \langle \psi,f \rangle = \int_{\mathbb{R}^n} d\boldsymbol{x}\ \overline{\psi(\boldsymbol{x})} f(\boldsymbol{x})
\end{equation}
Going forward, if we just write `state', we mean a generalised state.

\subsection{Dirac's braket notation in infinite dimensions}\label{subsec:braket_notation}

As in the finite-dimensional case, we can write states using Dirac's braket notation as kets $\ket{\psi}$.
Familiarity with the braket notation is maybe more of a hindrance than a help here!
For example, we can now write things like $\ket{\psi}(\phi)$ for $\phi \in \Phi$; this means we're applying the functional $\ket{\psi}$ to the Schwartz function $\phi$, giving us a complex number $\ket{\psi}(\phi)$.
This possibly looks quite weird if you're only used to finite-dimensional braket notation.

A bra $\bra{\psi}$ can then be defined to be the conjugate $\overline{\ket{\psi}}$.
That is, we define $\bra{\psi}(\phi) = \overline{\ket{\psi}}(\phi) = \overline{\ket{\psi}(\phi)}$.
The problems start when we want to define \textit{ketbras} $\ketbra{\psi}{\phi}$ and \textit{brakets} $\braket{\psi | \phi}$.
In analogy with the finite-dimensional case, we'd probably want $\ketbra{\psi}{\phi}$ to be something that maps a state to a state, and $\braket{\psi | \phi}$ to be a complex number.
We can achieve the former by defining $(\ketbra{\psi}{\phi})(\ket{\theta}) = \ket{\psi}\braket{\phi|\theta}$, so now if we can just find a way to interpret $\braket{\phi|\theta}$ as a complex number we're done.

Unfortunately, this isn't possible in full generality (or if it is, I can't find any references for it!).
But it is at least possible in all the cases we'll care about in this thesis.
Firstly, note that since $\Phi^* \cong \Phi^\times$ under the map $\ket{\psi} \mapsto \overline{\ket{\psi}} = \bra{\psi}$, we can freely identify bras and kets.
Also, we'll identify $\Phi$ and $\mathcal{H}$ with their copies in $\Phi^*$ and $\Phi^\times$ as in \Cref{eq:rigged_spaces_combined}.
So, when both $\ket{\psi} = f$ and $\ket{\phi} = g$ are in fact in the Hilbert space $\mathcal{H}$, then as in the finite-dimensional case we can define $\braket{\psi|\phi}$ to be the inner product $\langle \ket{\psi}, \ket{\phi} \rangle = \langle f, g \rangle = \int_{\mathbb{R}^n} d\boldsymbol{x} \overline{f(\boldsymbol{x})} g(\boldsymbol{x})$.
Next, when one of the two $\ket{\psi} = F$ is in $\Phi^* - \mathcal{H}^*$ but the other one $\ket{\phi} = \phi$ is strictly in $\Phi$, then we can define $\braket{\psi|\phi}$ to be the application of functional $\ket{\psi} = F$ to the Schwartz function $\ket{\phi} = \phi$; that is, $\braket{\psi|\phi} = \ket{\psi}(\ket{\phi}) = F(\phi)$.
If instead both $\ket{\psi}$ and $\ket{\phi}$ are in $\Phi^* - \mathcal{H}^*$, then unless we know something more about exactly what states $\ket{\psi}$ and $\ket{\phi}$ are, we're stuck.
We'll see examples of how to proceed in certain scenarios in the next subsections.
%For example, in the next subsection, we'll see how to proceed when $\ket{\psi}$ and $\ket{\phi}$ are \textit{generalised eigenstates} of the \textit{position and momentum operators}.

\subsection{Evolution}\label{subsec:qumode_evolution}

In the qubit case, we said pure states evolve via unitaries and measurements.
That's basically exactly the same here.
Indeed, physical states in $\HHH = \LLL^2{\RR^n}$ evolve under unitaries $\HHH \to \HHH$.
It's not entirely clear what it means for generalised states in $\Phi^\times$ to evolve -- evolution is about what physical processes can happen to states, and generalised states aren't `physical' in the first place.
Some sources declare that these states evolve under continuous linear maps $\Phi^\times \to \Phi^\times$, where as usual we won't go into what it means to be continuous.
What we will state is that any continuous linear map $\hat{f}: \HHH \to \HHH$ can be extended to one $\Phi^\times \to \Phi^\times$: for any $\ket{\psi} \in \Phi^\times$, define $\hat{f}(\ket{\psi})$ to be the functional that maps $\ket{\phi} \in \Phi$ to $\ket{\psi}(\hat{f}\ket{\phi})$.

\teague{Define measurements}

We'll also note that in the same way that the pure state qubit ZX-calculus wasn't restricted to representing unitaries and measurements, nor will the pure state Focked up ZX-calculus be.
So as in the qubit case, it won't be true that all ZX-diagrams correspond to physical processes that can be deterministically implemented.

\subsection{Fock states}

Time to meet our first continuous variable quantum states: the Fock states.
These are single-qumode states, and are not only physical (they're in $\HHH = \LLL^2(\RR)$) but are also in Schwartz space $\Phi \leq \HHH$.
We'll first just state the definition (which is pretty horrible looking) and then break it down.
For all $n \in \NN_0 = \{0, 1, \ldots\}$, we have:
\begin{equation}\label{eq:fock_state_defn}
    \ket{n}_F = \frac{1}{\sqrt{2^n n!}} \frac{1}{\pi^{1/4}} H_n(\frac{x}{\sqrt{2}}) e^{\frac{-x^2}{4}}
\end{equation}
The easy bit is $1/\sqrt{2^n n!} \cdot \pi^{1/4}$; this is just a real number, which gets (exponentially!) smaller as $n$ gets bigger.
Then we have $e^{-x^2/4}$ at the end -- this is sometimes called \textit{a Gaussian envelope}.
We can draw its graph:
\teague{TODO}
Finally, we have $H_n(\frac{x}{\sqrt{2}})$.
The $\frac{1}{\sqrt{2}}$ is just a rescaling of the interesting part $H_n(x)$, the $n$th Hermite polynomial.
These are defined recursively:
\begin{equation}\label{eq:hermite_polynomial_defn}
    \begin{aligned}
        H_0(x) &= 1 \\
        H_1(x) &= 2x \\
        H_n(x) &= 2x H_{n-1}(x) - 2(n-1)H_{n-2}(x) \text{ for $n \geq 2$}
    \end{aligned}
\end{equation}
As far as this thesis is concerned, I reckon the only thing you need to know about these (if anything!) is that $H_n$ is a degree $n$ polynomial.
Graphed, these look like:

So altogether, the Fock state $\ket{n}_F$ is essentially a degree $n$ polynomial $\mathbb{R} \to \mathbb{C}$, multiplied by a Gaussian envelope so that it flattens and decays to $0$ sufficiently quickly as $x$ tends to infinity, and multiplied by a prefactor so that TODO!
It's also smooth (infinitely differentiable).
Hence it's an element of not just the Hilbert space $\HHH = \LLL^2(\mathbb{R})$ but Schwartz space $\Phi \leq \HHH$.
Furthermore, it can be shown the Fock states form a basis for $\HHH$.
Hence we call $\{\ket{n}_F : n \in \NN_0\}$ the \textit{Fock basis}.

\subsection{Fourier transform}

Before introducing the next set of interesting states, we set out our Fourier transform conventions and state a few of its properties.
We will define the \textit{Fourier transform} $\mathcal{F}: \LLL^2(\RR^n) \to \LLL^2(\RR^n)$ to be:
\begin{equation}\label{eq:fourier_transform_defn}
    \mathcal{F}(f)(\bsym{p}) = \int_{\mathbb{R}^n} d\bsym{x}\ f(\bsym{x}) e^{-i2\pi\bsym{p} \cdot \bsym{x}}
\end{equation}
Then the inverse Fourier transform is:
\begin{equation}\label{eq:fourier_transform_inverse_defn}
    \mathcal{F}^{-1}(g)(\bsym{q}) = \int_{\mathbb{R}^n} d\bsym{x}\ g(\bsym{x}) e^{i2\pi\bsym{q} \cdot \bsym{x}}
\end{equation}
In the single-qumode case $n = 1$, it can be shown that the eigendecomposition of $\mathcal{F}$ in the Fock basis is:
\begin{equation}
    \mathcal{F} = \sum_{n \in \NN_0} i^n \ket{n}_F \bra{n}_F
\end{equation}
And the inverse Fourier transform has eigendecomposition:
\begin{equation}
    \mathcal{F}^{-1} = \sum_{n \in \NN_0} (-i)^n \ket{n}_F \bra{n}_F
\end{equation}
In the next chapter, we will make heavy reference to the function $\chi_a(x) \coloneqq e^{-i2\pi ax} : \mathbb{R} \to \mathbb{C}$, which winds around the real axis with a period of $1/a$.
This function, and its conjugation $\overline{\chi}_a(x) \coloneqq e^{i2\pi ax}$, satisfy the following two equations involving the Fourier transform:
\begin{equation}\label{eq:fourier_chi_relations}
    \mathcal{F}^{-1}(\mathcal{F}(f(x)) \cdot \chi_a) = f(x - a)
    \qqquad
    \mathcal{F}(\mathcal{F}^{-1}(f(x)) \cdot \overline{\chi}_a) = f(x - a)
\end{equation}

\subsection{Position and momentum eigenstates}

The next set of states we'll meet are the position and momentum eigenstates.
Unlike Fock states, these are not physical -- they're generalised states, i.e.\ continuous linear functionals $\Phi \to \CC$.
Indeed, they're generalised \textit{eigen}states.
An eigenvector of an operator $f: \Phi \to \Phi$ on Schwartz space with eigenvalue $\lambda \in \CC$ would be a function $\psi \in \Phi$ such that $f(\psi) = \lambda \psi$.
\teague{Is it really $\lambda \in \CC$? Or just in $\RR$?}
A \textit{(continuous) generalised eigenvector} of $f$ with eigenvalue $\lambda \in \CC$ is instead an element $F$ of the continuous dual space $\Phi^*$ such that $F(f(-)) = \lambda F(-)$.
That is, for all $\psi \in \Phi$, we must have $F(f(\psi)) = \lambda F(\psi)$.

So to define the position and momentum eigenstates, we must first define position and momentum operators.
These are operators on Schwartz space; they map functions in $\Phi$ to new functions in $\Phi$.
The position operator $\hat{q}$ is defined as $\hat{q}(\psi(x)) = x\psi(x)$.
The momentum operator $\hat{p}$ is defined as $\hat{p}(\psi(x)) = \frac{-i}{2\pi} \frac{d}{dx} \psi(x)$.
Note this is well-defined; Schwartz space operators are infinitely differentiable so $\frac{d}{dx} \psi(x)$ exists.
If you're new to CV quantum information, you might be wondering if it's annoying that the letter $p$ is used not for position but momentum -- you would be correct, I reckon.

A generalised eigenvector with eigenvalue $q \in \RR$ of the position operator $\hat{q}(\psi(x)) = x\psi(x)$ is a continuous linear functional $Q_q$ satisfying $Q_q(x(\psi(x))) = q Q_q(\psi(x))$.
The unique functional satisfying this condition is $Q_q(\phi(x)) = \phi(q)$, i.e.\ evaluation at $q$, since then $Q_q(x(\psi(x))) = q\psi(q) = q Q_q(\psi(x))$.
This functional is often called the \textit{Dirac distribution} and is denoted $\delta_q$.
Similarly, a generalised eigenvector with eigenvalue $p \in \RR$ of the momentum operator $\hat{p}(\psi(x)) = \frac{-i}{2\pi} \frac{d}{dx} \psi(x)$ is a continuous linear functional $P_p$ satisfying $P_p(\frac{-i}{2\pi} \frac{d}{dx} \psi(x)) = p P_p(\psi(x))$.
The unique such functional is $P_p(\phi(x)) = ?$, since then $P_p(\frac{-i}{2\pi} \frac{d}{dx} \psi(x)) = ? = p P_p(\psi(x))$.
In braket notation, we write position eigenstates $Q_q$ as $\ket{q}_Q$ and momentum eigenstates $P_p$ as $\ket{p}_P$.
We can once again show that they are bases for the space of all continuous linear functionals $\Phi^*$, so we refer to $\{\ket{q}_Q : q \in \mathbb{R}\}$ and $\{\ket{p}_P : p \in \mathbb{R}\}$ as the position and momentum bases respectively.
\teague{TODO verify everything!}

In the qubit case, the $X$ and $Z$ operators were related by conjugation by the Hadamard; $HXH^\dagger = Z$ and $HZH^\dagger = X$.
We will not prove it, but we have an analogous relationship here between the position and momentum operators via the Fourier transform;
namely, we have $\FFF \hat{q} \FFF^{-1} = -\hat{p}$ and $\FFF \hat{p} \FFF^{-1} = \hat{q}$.
This extends to the eigenstates too; since $\mathcal{F}$ is a continuous operator on $\Phi$, it can be extended to a continuous operator on $\Phi^\times$ as in \Cref{subsec:qumode_evolution}, and we get $\mathcal{F}\ket{q}_Q = \ket{q}_P$ and $\mathcal{F}\ket{p}_P = \ket{-p}_Q$. \teague{verify!!!}

\teague{Related by the Fourier transform}

\subsection{Fock, position and momentum: brakets}

In \Cref{subsec:braket_notation}, we said bras were just conjugates of kets, and ketbras can be defined in terms of brakets.
So to be able to use all four of these objects, it just remains to define brakets $\presub{A}{\braket{a | b}}_B$, where $A, B \in \{F, Q, P\}$, and $a$ and $b$ in $\mathbb{N}$ or $\mathbb{R}$ accordingly.
For Fock states, this is no problem, because they're in Schwartz space.
If both the bra and ket are Fock states, the braket ${}_F\braket{n | m}_F$ is defined using the inner product on $\HHH = \LLL^2(\mathbb{R})$; this is an ugly calculation but the end result is that ${}_F\braket{n | m}_F = \delta_{n,m}$, the usual Kronecker delta (not to be confused with the Dirac distribution!).
If just one is a Fock state, the brakets ${}_Q\braket{q | n}_F$ and ${}_P\braket{p | n}_F$ are defined as the application of the functional ${}_Q\bra{q}$ or ${}_P\bra{p}$ to the Schwartz function $\ket{n}_N$.
The result of this is that ${}_Q\braket{q | n}_F = ?$ and ${}_P\braket{p | n}_F = ?$ \teague{TODO -- if they're nice. Else just leave it as is.}
If neither are Fock states, then we have three cases.
First, both could be position eigenstates, so we have ${}_Q\braket{q | q'}_Q$.
TODO.
Alternatively, both could be momentum eigenstates, i.e.\ ${}_P\braket{p | p'}_P$.
TODO.
Else, we have one of each: ${}_Q\braket{q | p}_P$.
TODO.

\subsection{Phase space}\label{subsec:phase_space}

Rather than view a physical state $f \in \LLL^2(\RR^n)$ as a function $\mathbb{R}^n \to \mathbb{C}$, it's common to transform it into its \textit{Wigner function} $W_f : \RR^{2n} \to \RR$.
In this context, $\mathbb{R}^{2n}$ is called \textit{phase space}.
Throughout this section, we will often write vectors in phase space in block form $\tbinom{\bsym{v}}{\bsym{u}}$.
We say vectors are in \textit{$qqpp$ order}; suppose we were to reorder $\tbinom{\bsym{v}}{\bsym{u}}$ to $(v_1, u_1, \ldots, v_n, u_n)$, then this would be in \textit{$qpqp$ order}.
Given $f \in \LLL^2(\RR^n)$, its Wigner function $W_f$ is defined as:
\begin{equation}\label{eq:wigner_function_defn}
    W_f(\tbinom{\bsym{q}}{\bsym{p}}) = \int_{\mathbb{R}^n} d\bsym{x} f(\bsym{q}+\frac{\bsym{x}}{2}) \overline{f(\bsym{q}-\frac{\bsym{x}}{2})} e^{-i2\pi \bsym{p} \cdot \bsym{x}}
\end{equation}
We can break this definition down; first, let's define the intermediate function $\RR^n \to \CC$:
\begin{equation}
    f_{\bsym{q}}(\bsym{x}) = f(\bsym{q}+\frac{\bsym{x}}{2}) \overline{f(\bsym{q}-\frac{\bsym{x}}{2})}
\end{equation}
So the Wigner function is just the Fourier transform of the intermediate function:
\begin{equation}
    W_f(\tbinom{\bsym{q}}{\bsym{p}}) = \mathcal{F}(f_{\bsym{q}})(\bsym{p})
\end{equation}

We can just about extend all this to generalised states too -- that is, to continuous linear functionals.
Keeping to the theme of this section, we won't go into the maths in detail (we'd need to justify what it means for a continuous linear functional to appear in an integral, for example), but will just give the result for the position and momentum eigenstates $Q_q$ and $P_p$:
\begin{equation}
    W_{Q_q}(x, y) = \delta_{x - q}
    \qquad \text{and} \qquad
    W_{P_p}(x, y) = \delta_{y - p}
\end{equation}
Here $\delta_{a}$ is the \textit{Dirac distribution} mentioned above, that takes a Schwartz function $\psi$ and evaluates it at $a$, i.e.\ $\delta_a(\psi) = \psi(a)$.
\teague{Phase space is nice because TODO! MAkes symmetry between position and momentum obvious. Also because it makes clear why displacements are called displacements}

\subsection{Displacements}

Displacement operators can be viewed as the continuous variable analogue of the Pauli operators.
Writing $\HHH = \LLL^2(\RR^n)$, a displacement operator is any operator of the form $e^{i\theta} D(\bsym{v}, \bsym{u}): \HHH \to \HHH$ for $\bsym{v}, \bsym{u} \in \RR^{n}$, where:
\begin{equation}\label{eq:displacement_defn}
    D(\tbinom{\bsym{v}}{\bsym{u}})
    \coloneqq e^{-i2\pi \tbinom{\bsym{v}}{\bsym{u}} \odot \tbinom{\bsym{\hat{q}}}{\bsym{\hat{p}}}}
    = e^{-i2\pi (\bsym{v} \cdot \bsym{\hat{p}} - \bsym{u} \cdot \bsym{\hat{q}})}
\end{equation}
Equivalently, we can define it by what it does to an element $f \in \HHH$:
\begin{equation}\label{eq:displacement_defn_action}
    D(\tbinom{\bsym{v}}{\bsym{u}})(f(\bsym{x})) = e^{-i\pi \bsym{v} \cdot \bsym{u}} e^{i2\pi \bsym{x} \cdot \bsym{u}} f(\bsym{x} - \bsym{v})
\end{equation}
But perhaps the more intuitive definition, and the reason for its name, is via its action on position and momentum operators under conjugation, or equivalently its action on Wigner functions.
The former is:
\begin{equation}
    D(\tbinom{\bsym{v}}{\bsym{u}}) \hat{\bsym{q}} D(\tbinom{\bsym{v}}{\bsym{u}})^\dagger = \hat{q} - \bsym{v}
    \qqquad
    D(\tbinom{\bsym{v}}{\bsym{u}}) \hat{\bsym{p}} D(\tbinom{\bsym{v}}{\bsym{u}})^\dagger = \hat{p} - \bsym{u}
\end{equation}
While the latter is:
\begin{equation}
    W_{D(\tbinom{\bsym{v}}{\bsym{u}})(f)}(\bsym{q}, \bsym{p}) = W_f(\bsym{q} - \bsym{v}, \bsym{p} - \bsym{u})
\end{equation}
Displacement operators obey the following (anti-)commutation relation:
\begin{equation}\label{eq:displacement_commutation}
    D(\tbinom{\bsym{v}}{\bsym{u}}) D(\tbinom{\bsym{v}'}{\bsym{u}'}) = e^{-i 2 \pi \tbinom{\bsym{v}}{\bsym{u}} \odot \tbinom{\bsym{v}'}{\bsym{u}'}} D(\tbinom{\bsym{v}'}{\bsym{u}'}) D(\tbinom{\bsym{v}}{\bsym{u}})
\end{equation}
And the product of any two displacement operators is again a displacement, with vectors $\tbinom{\bsym{v}}{\bsym{u}}$ and $\tbinom{\bsym{v}'}{\bsym{u}'}$ adding together:
\begin{equation}\label{eq:displacement_product_is_displacement}
    D(\tbinom{\bsym{v}}{\bsym{u}}) D(\tbinom{\bsym{v}'}{\bsym{u}'}) = e^{-i \pi \tbinom{\bsym{v}}{\bsym{u}} \odot \tbinom{\bsym{v}'}{\bsym{u}'}}  D(\tbinom{\bsym{v}}{\bsym{u}} + \tbinom{\bsym{v}'}{\bsym{u}'})
\end{equation}

\subsection{Gaussian unitaries}

TODO!